\section{Practical Considerations}\label{sec:practical}

\subsection{Formatting the Conditioning Context}

Computing $\theta(r_k \mid r_{<k}, p)$ requires feeding $\theta$ a context containing the prompt and previous responses.
The formatting choice affects results.
We recommend using a base (non-instruction-tuned) model as $\theta$ and formatting as:

\begin{verbatim}
    [prompt p]

    Response A: [r_1]

    Response B: [r_2]

    Response C: [r_3]
\end{verbatim}

This leverages the base model's in-context learning to treat the responses as parallel completions.
An instruction-tuned model might instead interpret previous responses as conversation turns, confounding the measurement.

\paragraph{Base vs.\ instruction-tuned models.} We use a base (non-instruction-tuned) model as $\theta$ for two reasons.
First, an instruction-tuned model has had its output distribution shaped by RLHF or similar procedures, so using one as $\theta$ reintroduces the kind of distributional bias we seek to avoid.
Second, an instruction-tuned $\theta$ may assign systematically different probabilities to text that follows or violates its alignment training, introducing a confound between coherence-as-fluency and coherence-as-alignment that the metric should not conflate.
Modern base models already exhibit strong in-context learning from pretraining alone, so this choice does not sacrifice ICL capability.
Optimizing the prompt format to maximally elicit in-context learning from $\theta$ is an interesting direction but out of scope for this paper; we use the neutral ``Response A/B/C'' template throughout.

\subsection{Computational Cost}\label{sec:cost}

Because $\theta$ is a causal language model, the full $a_k$ curve can be computed from a \textbf{single forward pass}.
Concatenate the prompt and all responses into one sequence:

\begin{verbatim}
    [prompt p] Response A: [r_1] Response B: [r_2] ... Response N: [r_n]
\end{verbatim}

The forward pass produces the log-probability of every token conditioned on all preceding tokens.
The tokens belonging to $r_k$ are conditioned on exactly $p, r_1, \ldots, r_{k-1}$ (plus formatting), which is the conditioning required for $a_k$.
Partitioning the output log-probabilities by response boundary and summing within each partition yields all $n$ values of $a_k$ simultaneously.
The cost is a single forward pass over the full concatenation, with FLOPs scaling as $O((|p| + n\bar{L}_{\mathrm{tok}})^2)$ from causal attention, where $\bar{L}_{\mathrm{tok}}$ is the average response length in tokens.

The coherence term $C$ and spread $\sigma_\ell$ additionally require the \emph{unconditional} per-byte cross-entropies $h_\theta(r_i \mid p)$ for each response.
These cannot be extracted from the concatenated pass, since in that pass $r_i$ for $i > 1$ is conditioned on all prior responses.
Computing them requires $n$ independent forward passes, each over the short context $(p, r_i)$.
These passes are embarrassingly parallel and batchable.

In summary:
\begin{itemize}[nosep]
    \item $a_k$ curve: 1 forward pass per permutation (long context).
    \item $C$ and $\sigma_\ell$: $n$ forward passes (short contexts, batchable).
\end{itemize}

\subsection{Dependence on Sample Ordering}\label{sec:ordering}

The $a_k$ values depend on the ordering of $\{r_i\}$.
Since raw $a_k$ is in total bits, responses of different lengths produce different values even at the same per-byte rate, making a single ordering's curve jagged.
Averaging over random permutations removes this length effect: each position averages over all responses, so the curve reflects only how $\theta$'s predictions improve with more context.
We recommend averaging over a moderately large number of random permutations; Section~\ref{sec:practical-findings} provides empirical evidence that $\geq 50$ permutations are needed for reliable rankings.

Per-response normalization by byte count must be performed \emph{before} averaging across permutations, since the byte count depends on which response lands at each position.
Concretely, the per-byte curve is a mean of per-permutation rates,
\begin{equation}\label{eq:akbar-mor}
    \bar{a}_k = \frac{1}{|\Sigma|}\sum_{\sigma \in \Sigma} \frac{a_k^{\sigma}}{\|r_{\sigma(k)}\|},
\end{equation}
not a ratio of permutation-averaged bits to permutation-averaged byte counts (which would degenerate, after enough permutations, into the total-bits curve rescaled by the mean response length).

\subsection{Choice of $n$}

The diversity score depends on $n$ through the estimate $a_n \approx a_\infty$.
We recommend choosing $n$ large enough that $a_n$ has approximately stabilized (i.e., the curve has reached its floor), and fixing $n$ across all comparisons for consistency.
A practical choice is $n = 10$--20, which suffices for policies where $\theta$'s learning stabilizes within this range.

\paragraph{Detectable mode weight.} Modes with weight $w \lesssim 1/n$ will rarely appear in the sample, and even when they do, $\theta$ has insufficient examples to learn the pattern.
As a rough guide, modes with weight below $1/n$ are invisible to the metric.
For $n = 15$, this means modes contributing less than $\sim$7\% of the probability mass are unlikely to be captured.
If rare modes are of interest (e.g., when evaluating whether an intervention eliminates low-frequency behaviors), $n$ should be increased accordingly, subject to the context-length constraints discussed above.

%%%%%%%%%%%%%%%%%%%%%%%%%%%%%%%%%%%%%%%%%%%%%%%%%%%%%%%%%%%%%%%%%%%%%%%%
%% SECTION 8: EXPERIMENTS
%%%%%%%%%%%%%%%%%%%%%%%%%%%%%%%%%%%%%%%%%%%%%%%%%%%%%%%%%%%%%%%%%%%%%%%%
